\par\Needspace{6\baselineskip}\medskip
\hypertarget{native-115-article086}{}\pdfbookmark[2]{86.}{article-086}
\begin{center}86.\end{center}
Designando radium vectorem indefinitum seu variabilem anomaliae verae $v-\Pi$ respondentem per $\rho$, erit area sectoris a corpore coelesti intra tempus $t$ descripti $=\tfrac12\int\rho\rho\,\mathrm{d}v$, hoc integrali a $v=N$ usque ad $v=N'$ extenso, adeoque, accipiendo $k$ in significatione art. 6, $kt\BareRoot{p}=\int\rho\rho\,\mathrm{d}v$. Iam constat, per formulas a \textsc{Cotesio} evolutas, si $\varphi x$ exprimat functionem quamcunque ipsius $x$, valorem continuo magis approximatum integralis $\int\varphi x\mathbin{.}\mathrm{d}x$ ab $x=u$ usque ad $x=u+\Delta$ extensi exhiberi per formulas
\[
\begin{aligned}
&\tfrac12\Delta\{\varphi u+\varphi(u+\Delta)\}\\[4pt]
&\tfrac16\Delta\{\varphi u+4\varphi(u+\tfrac12\Delta)+\varphi(u+\Delta)\}\\[4pt]
&\tfrac18\Delta\{\varphi u+3\varphi(u+\tfrac13\Delta)+3\varphi(u+\tfrac23\Delta)+\varphi(u+\Delta)\}
\end{aligned}
\]
\centerline{etc.:}
\noindent ad institutum nostrum apud duas formulas primas subsistere sufficiet.

Per formulam itaque primam in problemate nostro habemus $\int\rho\rho\,\mathrm{d}v=\tfrac12\Delta(rr+r'r')=\dfrac{\Delta rr'}{\cos2\omega}$, si statuitur $\dfrac{r'}r=\operatorname{tang}(45^0+\omega)$. Quamobrem valor approximatus primus ipsius $\BareRoot{p}$ erit $=\dfrac{\Delta rr'}{kt\cos2\omega}$, quem statuemus $=3a$.

Per formulam secundam habemus exactius $\int\rho\rho\,\mathrm{d}v=\tfrac16\Delta(rr+r'r'+4RR)$, designante $R$ radium vectorem anomaliae intermediae $\tfrac12N+\tfrac12N'-\Pi$ respondentem. Iam exprimendo $p$ per $r$, $R$, $r'$, $N$, $N+\tfrac12\Delta$, $N+\Delta$ ad normam formulae in art. 82 traditae, invenimus
\[
p=\frac{4\sin\tfrac14\Delta^2\sin\tfrac12\Delta}{\left(\dfrac1r+\dfrac1{r'}\right)\sin\tfrac12\Delta-\dfrac1R\sin\Delta}\mathpunct{\raisebox{.45ex}{,}}
\]
\SourceLeaf{116}atque hinc
\[
\begin{aligned}
\frac{\cos\tfrac12\Delta}{R}&=\tfrac12\left(\frac1r+\frac1{r'}\right)-\frac{2\sin\tfrac14\Delta^2}{p}\\[4pt]
&=\frac{\cos\omega}{\BareRoot{(rr'\cos2\omega)}}-\frac{2\sin\tfrac14\Delta^2}{p}\mathpunct{\cdot}
\end{aligned}
\]
Statuendo itaque
\[
\frac{2\sin\tfrac14\Delta^2\BareRoot{(rr'\cos2\omega)}}{\cos\omega}=\delta,
\]
fit
\[
R=\frac{\cos\tfrac12\Delta\mathbin{.}\BareRoot{(rr'\cos2\omega)}}{\cos\omega\mathbin{.}\left(1-\dfrac{\delta}{p}\right)}\mathpunct{\raisebox{.45ex}{,}}
\]
unde valor approximatus secundus ipsius $\BareRoot{p}$ elicitur
\[
\BareRoot{p}=a+\frac{2a\cos\tfrac12\Delta^2\cos2\omega^2}{\cos\omega^2\left(1-\dfrac{\delta}{p}\right)^2}=a+\frac{\epsilon}{\left(1-\dfrac{\delta}{p}\right)^2}\mathpunct{\raisebox{.45ex}{,}}
\]
si statuitur $2a\left(\dfrac{\cos\tfrac12\Delta\cos2\omega}{\cos\omega}\right)^2=\epsilon$. Scribendo itaque $\pi$ pro $\BareRoot{p}$, determinabitur $\pi$ per aequationem $(\pi-a)\left(1-\dfrac{\delta}{\pi\pi}\right)^2=\epsilon$, quae rite evoluta ad quintum gradum ascenderet. Statuamus $\pi=q+\mu$, ita ut sit $q$ valor approximatus ipsius $\pi$, atque $\mu$ quantitas perexigua, cuius quadrata altioresque potestates negligere liceat:

Qua substitutione prodit
\[
\begin{aligned}
&(q-a)\left(1-\frac{\delta}{qq}\right)^2\\[4pt]
&\quad+\mu\left\{\left(1-\frac{\delta}{qq}\right)^2+\frac{4\delta(q-a)}{q^3}\left(1-\frac{\delta}{qq}\right)\right\}=\epsilon,
\end{aligned}
\]
sive
\[
\mu=\frac{\epsilon q^5-(qq-aq)(qq-\delta)^2}{(qq-\delta)(q^3+3\delta q-4a\delta)}\mathpunct{\raisebox{.45ex}{,}}
\]
adeoque
\[
\pi=\frac{\epsilon q^5+(qq-\delta)(aqq+4\delta q-5a\delta)q}{(qq-\delta)(q^3+3\delta q-4a\delta)}\mathpunct{\cdot}
\]
Iam in problemate nostro habemus valorem approximatum ipsius $\pi$, puta $=3a$, quo in formula praecedente pro $q$ substituto, prodit valor correctus
\[
\pi=\frac{243a^4\epsilon+3a(9aa-\delta)(9aa+7\delta)}{(9aa-\delta)(27aa+5\delta)}\mathpunct{\cdot}
\]
Statuendo itaque $\dfrac{\delta}{27aa}=\beta$, $\dfrac{\epsilon}{(1-3\beta)a}=\gamma$, formula induit formam hancce $\pi=\dfrac{a(1+\gamma+21\beta)}{1+5\beta}$, omnesque operationes ad problematis solutionem necessariae in his quinque formulis continentur:
\[
\begin{aligned}
\mathrm{I.}\qquad&\frac{r'}{r}=\operatorname{tang}(45^0+\omega)\\[6pt]
\mathrm{II.}\qquad&\frac{\Delta rr'}{3kt\cos2\omega}=a
\end{aligned}
\]
\SourceLeaf{117}
\[
\begin{aligned}
\mathrm{III.}\qquad&\frac{2\sin\tfrac14\Delta^2\BareRoot{(rr'\cos2\omega)}}{27aa\cos\omega}=\beta\\[6pt]
\mathrm{IV.}\qquad&\frac{2\cos\tfrac12\Delta^2\cos2\omega^2}{(1-3\beta)\cos\omega^2}=\gamma\\[6pt]
\mathrm{V.}\qquad&\frac{a(1+\gamma+21\beta)}{1+5\beta}=\BareRoot{p}.
\end{aligned}
\]

Si quid a praecisione harum formularum remittere placet, expressiones adhuc simpliciores evolvere licebit. Scilicet faciendo $\cos\omega$ et $\cos2\omega=1$ et evolvendo valorem ipsius $\BareRoot{p}$ in seriem secundum potestates ipsius $\Delta$ progredientem, prodit neglectis biquadratis altioribusque potestatibus
\[
\BareRoot{p}=a\left(3-\tfrac12\Delta\Delta+\frac{\Delta\Delta\BareRoot{rr'}}{18aa}\right),
\]
ubi $\Delta$ in partibus radii exprimendus est. Quare faciendo $\dfrac{\Delta rr'}{kt}=\BareRoot{p'}$, habetur
\[
\mathrm{VI.}\qquad p=p'\left(1-\tfrac13\Delta\Delta+\frac{\Delta\Delta\BareRoot{rr'}}{3p'}\right).
\]
Simili modo explicando $\BareRoot{p}$ in seriem secundum potestates ipsius $\sin\Delta$ progredientem emergit, posito $\dfrac{rr'\sin\Delta}{kt}=\BareRoot{p''}$,
\[
\mathrm{VII.}\qquad\BareRoot{p}=\left(1+\frac{\sin\Delta^2\BareRoot{rr'}}{6p''}\right)\BareRoot{p''}
\]
sive
\[
\mathrm{VIII.}\qquad p=p''+\tfrac13\sin\Delta^2\BareRoot{rr'}.
\]
Formulae VII et VIII conveniunt cum iis, quas ill. \textsc{Euler} tradidit in \textit{Theoria motus planetarum et cometarum} [p. 56], formula VI autem cum ea, quae in usum vocata est in \textit{Recherches et calculs sur la vraie orbite elliptique de la comète de 1769}, p. 80.
